\documentclass[11pt]{article}
\usepackage[margin=1in]{geometry}
\usepackage{amsmath,amssymb,graphicx,booktabs,hyperref,xcolor}
\title{Computational Replication Report:\\ Zhang \emph{et al.} (2025),\\ \emph{Strain-induced nonrelativistic altermagnetic spin splitting effect} (arXiv:2503.17916)}
\author{Ollie (OpenClaw @ Argonne) --- TEXTURES-100 Wave 4}
\date{2026-07-19}
\begin{document}
\maketitle

\section{Paper summary}
Using first-principles DFT ($+U$) with Wannier interpolation and a constant-$\Gamma$ Kubo
linear-response evaluation, the paper proposes \emph{strain engineering} to induce
altermagnetism in rutile OsO$_2$ (and RuO$_2$). Equibiaxial tensile strain $E_{ts}\gtrsim 2\%$
drives a nonmagnetic$\to$altermagnetic transition; the Os moment and the nonrelativistic band
spin splitting both rise then fall with strain (dome, peak $\sim$4--5\%), and the altermagnetic
spin-splitting effect (ASSE) yields a spin--charge conversion ratio $\theta_{AS}\!\sim\!7\%$,
persisting without SOC (nonrelativistic origin). Published as Phys.\ Rev.\ B \textbf{112}, 024415.

\section{Scope statement}
Full DFT+Wannier+Kubo is out of scope on CPU-only hardware. We replicate (i) the paper's OWN
quantitative correlation from Table~I, (ii) the non-monotonic strain dome via a minimal Stoner
mean field, and (iii) the SOC-free (nonrelativistic) moment$\to$splitting mechanism in a $d$-wave
tight-binding model. Absolute meV splittings and $\theta_{AS}$ are DFT-only and reported, not recomputed.

\section{Claims addressed}
\begin{table}[h]\centering\small
\begin{tabular}{clccc}
\toprule
ID & Claim & Type & Testable (CPU)? & Tested? \\
\midrule
C1 & Spin splitting positively correlates with strain-induced moment & data corr & yes & yes \\
C2 & Non-monotonic dome; AM onsets $>\!2\%$, peaks $\sim$4--5\% & mechanism & yes & yes \\
C3 & Splitting persists WITHOUT SOC (nonrelativistic origin) & model & yes & yes \\
C4 & $\theta_{AS}\!\sim\!7\%$ ASSE conversion ratio & Kubo/DFT & no & \textcolor{red}{reported only} \\
\bottomrule
\end{tabular}
\end{table}

\section{Paper Table I (reproduced input)}
\begin{table}[h]\centering\small
\begin{tabular}{lccccccc}
\toprule
$E_{ts}$ (\%) & 0 & 1 & 2 & 3 & 4 & 5 & 6 \\
Os moment ($\mu_B$) & 0.000 & 0.002 & 0.028 & 0.349 & 0.468 & 0.500 & 0.150 \\
$|\Delta|_{\max}$ (meV) & 0 & 62.8 & 209.2 & 256.3 & 275.5 & 110.0 & 0.11 \\
\bottomrule
\end{tabular}
\end{table}

\section{Results vs.\ paper}
\begin{table}[h]\centering\small
\begin{tabular}{lccc}
\toprule
Quantity & Paper & This work & Verdict \\
\midrule
moment--splitting Spearman $\rho$ & positive & $0.643$ & \checkmark \\
Pearson (rising branch $E_{ts}\!\le\!4$) & positive & $0.807$ & \checkmark (sub-linear) \\
AM onset strain & $>\!\sim\!2\%$ & $\sim$2.5\% (model) & \checkmark \\
dome peak & $\sim$4--5\% & $\sim$4.3\% (model); Table I 4--5\% & \checkmark \\
SOC-free splitting vs $m$ slope & linear, nonzero & $8.00$, residual $10^{-15}$ & \checkmark \\
$\theta_{AS}$ & $\sim7\%$ & not recomputed & out of scope \\
\bottomrule
\end{tabular}
\end{table}

\section{Per-claim what worked / what didn't}
\textbf{C1 (partial/worked).} Positive association confirmed (Spearman $0.64$), and near-linear on
the rising branch ($0.81$); a single Pearson over all points ($0.42$) is lowered by the $E_{ts}=$5--6\%
dome collapse --- an honest nonlinearity (splitting saturates faster than moment, Open Q1).
\textbf{C2 (worked).} The Stoner dome reproduces onset $>\!2\%$ and peak $\sim$4--5\%, matching Table~I
(moment peak 5\%, splitting peak 4\%). \textbf{C3 (worked).} The SOC-free $d$-wave model gives a
nonzero splitting linear in the moment ($8\lambda m$), realizing the nonrelativistic origin.
\textbf{C4 (out of scope).} $\theta_{AS}\!\sim\!7\%$ needs the material Kubo/Wannier calculation.

\section{Critique of the replication}
This is a \emph{mechanism + data-correlation} replication, NOT an independent DFT reproduction ---
the absolute meV numbers come from the paper's Table~I, which we analyze rather than regenerate.
The Stoner dome and the linear $t_{\rm AM}\!\propto\!m$ TB are deliberately minimal: they capture the
QUALITATIVE physics (onset, dome, nonrelativistic persistence) but not the quantitative saturation
of the splitting nor $\theta_{AS}$. The LLM-judge scored PARTIAL (coverage 7, agreement 6),
which we adopt: the mechanism is confirmed, the DFT/transport headline is method-limited.
This is a canonical reduced-model-vs-full-DFT PARTIAL for this campaign.

\section{Verdict}
\textbf{PARTIAL} --- mechanism and the paper's own moment--splitting correlation + non-monotonic
strain dome + nonrelativistic (SOC-free) origin reproduced; absolute meV splittings and
$\theta_{AS}\!\sim\!7\%$ require DFT+Wannier+Kubo (out of scope, documented as method-limited).

\section*{Open Questions}
See \texttt{report/open\_questions.json}: Q1 sub-linear moment$\to$splitting saturation; Q2 6\%
collapse mechanism (Lifshitz vs Stoner); Q3 $\theta_{AS}$ sensitivity to $\Gamma$ and $U$; Q4 RuO$_2$
magnetism controversy; Q5 $d$-wave-vs-FM discrimination of the strained state.

\end{document}
